% ============================================================
%  HEATR Simulation Methods
%  Discretization, Thermal Model, EQS Solver, and Densification
% ============================================================
\documentclass[11pt, letterpaper]{article}

\usepackage[margin=1.1in]{geometry}
\usepackage{amsmath, amssymb}
\usepackage{bm}
\usepackage{booktabs}
\usepackage{xcolor}
\usepackage{hyperref}
\usepackage{microtype}
\usepackage{parskip}
\usepackage{titlesec}
\usepackage{enumitem}
\usepackage{nicefrac}

\definecolor{heatrblue}{HTML}{3B82F6}
\definecolor{ruleclr}{HTML}{CBD5E1}

\hypersetup{colorlinks=true, linkcolor=heatrblue, citecolor=heatrblue, urlcolor=heatrblue}

\titleformat{\section}{\large\bfseries}{}{0em}{}[\color{ruleclr}\hrule]
\titleformat{\subsection}{\normalsize\bfseries\color{heatrblue}}{}{0em}{}
\titlespacing{\section}{0pt}{1.4em}{0.5em}
\titlespacing{\subsection}{0pt}{1.0em}{0.3em}

% Shorthand
\newcommand{\Qrf}{Q_{\mathrm{rf}}}
\newcommand{\rhorel}{\rho_{\mathrm{rel}}}
\newcommand{\cpeff}{c_{p,\mathrm{eff}}}
\newcommand{\phimelt}{\phi_{\mathrm{melt}}}
\newcommand{\partmask}{\Omega_{\mathrm{part}}}

% ============================================================
\begin{document}

\begin{center}
  {\LARGE \textbf{HEATR: Simulation Methods Reference}}\\[0.4em]
  {\large High-frequency Electrothermal Additive Thermal Resolver}\\[0.6em]
  {\small\color{gray}Derived from \texttt{rfam\_eqs\_coupled.py} --- April 2026}
\end{center}

\vspace{0.3em}\noindent\color{ruleclr}\hrule\vspace{1em}

\noindent
HEATR is a 2D coupled electroquasistatic--thermal--densification solver for
Radio-Frequency Additive Manufacturing (RFAM) of polymer powder beds.
This document describes the mathematical formulation, numerical discretization,
material models, and time-integration strategy implemented in the solver.

% ============================================================
\section{Spatial Domain and Discretization}

\subsection{Computational Grid}

The simulation domain is a 2D cross-section of the RF chamber in the $xy$-plane.
Variation in the depth direction ($z$) is not directly resolved; instead it is
accounted for by an analytical correction term (Section~\ref{sec:depth}).

The domain is discretized onto a uniform rectangular grid of
$N_x \times N_y$ nodes (default: $120 \times 120$):
%
\begin{equation}
  x_i = -\tfrac{L_x}{2} + \frac{i\,L_x}{N_x - 1},\quad i = 0,\ldots,N_x-1,
  \qquad
  y_j = -\tfrac{L_y}{2} + \frac{j\,L_y}{N_y - 1},\quad j = 0,\ldots,N_y-1,
\end{equation}
%
where $L_x \times L_y$ is the chamber width and height (default: $60\times60$~mm).
Grid spacings are $\Delta x = L_x/(N_x-1)$ and $\Delta y = L_y/(N_y-1)$,
giving $\approx 0.5$~mm per cell at the defaults.

\subsection{Part Rasterization}

The part geometry is defined as a closed polygon and rasterized onto the grid
using a \emph{subpixel fill-fraction} scheme.
Each cell on or near the part boundary is assigned a fill fraction
$f_{ij}\in[0,1]$ proportional to the fraction of the cell interior that lies
inside the polygon.
This anti-aliasing approach avoids the staircase artifacts that arise from
hard binary rasterization, and the fill fraction is used directly in the
material property interpolation and the EQS conductivity field.

% ============================================================
\section{Electroquasistatic (EQS) Solver}

\subsection{Governing Equation}

At RF frequencies ($\sim$27~MHz), inductive effects and wave propagation
are negligible at the chamber scale.
The electric field is governed by the 2D electroquasistatic Laplace equation:
%
\begin{equation}
  \label{eq:eqs}
  \nabla\cdot\bigl[\gamma(\bm{r})\,\nabla V(\bm{r})\bigr] = 0,
\end{equation}
%
where $V$ is the complex electric potential and $\gamma$ is the
\emph{complex admittivity}:
%
\begin{equation}
  \gamma = \sigma + j\omega\varepsilon_0\varepsilon_r.
\end{equation}
%
Here $\sigma$~[S/m] is electrical conductivity, $\varepsilon_r$ is the relative
permittivity, and $\omega = 2\pi f$ is the angular RF frequency.
Both $\sigma$ and $\varepsilon_r$ vary spatially depending on whether a cell
is inside the doped part, in the surrounding virgin powder, or at a mixed
(subpixel) boundary.
The admittivity field is held constant within each $\Delta t$ interval,
which is justified because the RF period ($\sim$37~ns) is orders of magnitude
shorter than the thermal timescale ($\sim$1~s).

\subsection{Discrete Formulation}

Equation~\eqref{eq:eqs} is discretized on the rectangular grid by finite
differences.
At each interior node $(i,j)$ that is not an electrode, the divergence is
replaced by a flux balance over the four face-adjacent neighbours.
At each cell face, the admittivity is \emph{harmonically averaged}:
%
\begin{equation}
  \gamma_{\mathrm{face}} = \frac{2\,\gamma_{ij}\,\gamma_{\mathrm{nb}}}
                                {\gamma_{ij} + \gamma_{\mathrm{nb}}},
\end{equation}
%
where $\gamma_{\mathrm{nb}}$ is the admittivity of the adjacent cell.
Harmonic averaging is used rather than arithmetic averaging because it correctly
limits conductance when one of the two cells is nearly insulating.
Arithmetic averaging overestimates interfacial flux across the high-contrast
doped-part/virgin-powder boundary, which would produce spurious
corner-heating artifacts.

This yields the sparse complex linear system
$\mathbf{A}\bm{v} = \bm{b}$,
where $\mathbf{A}\in\mathbb{C}^{N\times N}$ ($N = N_x N_y$) encodes the
interior difference stencils and $\bm{b}$ encodes the Dirichlet electrode
conditions.
The system is solved \emph{exactly} by direct sparse LU factorization
(\texttt{scipy.sparse.linalg.spsolve}) with a small diagonal regularization
$\varepsilon_{\mathrm{reg}}=10^{-18}$ to improve conditioning in nearly-insulating
regions.

\subsection{Boundary Conditions}

\begin{itemize}[leftmargin=2em, itemsep=2pt]
  \item \textbf{Top and bottom edges (electrodes) --- Dirichlet.}
    In \emph{centered} mode: $V_\mathrm{hi}=+V_0/2$, $V_\mathrm{lo}=-V_0/2$.
    In \emph{grounded} mode: $V_\mathrm{hi}=V_0$, $V_\mathrm{lo}=0$.
  \item \textbf{Left and right edges --- Neumann} (zero normal flux;
    electrically insulating walls).
\end{itemize}

\subsection{Electric Field and Power Density}

The electric field components are recovered by central differences on the
solved potential:
%
\begin{equation}
  E_x = -\frac{\partial V}{\partial x},\qquad E_y = -\frac{\partial V}{\partial y}.
\end{equation}
%
The time-averaged volumetric RF power density (resistive heating source) is:
%
\begin{equation}
  \label{eq:qrf}
  \Qrf = \tfrac{1}{2}\,\mathrm{PF}\cdot\mathrm{Re}\!\left[\gamma\,|\bm{E}|^2\right]
       = \tfrac{1}{2}\,\mathrm{PF}\cdot\sigma\,|\bm{E}|^2
  \quad[\mathrm{W/m^3}],
\end{equation}
%
where PF is a power-factor correction (default 1.0).
$\Qrf$ is clipped to zero to suppress sub-zero noise from discretization.

When generator-power enforcement is active, $\Qrf$ is globally rescaled so that
%
\begin{equation}
  \sum_{ij}\Qrf\,\Delta x\,\Delta y\cdot d_{\mathrm{eff}} = P_{\mathrm{gen}}\,\eta,
\end{equation}
%
where $d_{\mathrm{eff}}$ is the effective chamber depth and $\eta$ is the RF
transfer efficiency.

% ============================================================
\section{Thermal Solver}
\label{sec:thermal}

\subsection{Governing Equation}

Inside the part domain $\partmask$, the temperature field evolves according to
the heat equation with volumetric source and surface sinks:
%
\begin{equation}
  \label{eq:heat}
  \rho\,\cpeff\,\frac{\partial T}{\partial t}
  = \nabla\cdot\bigl[k(\bm{r},T)\,\nabla T\bigr]
    + \Qrf
    - q_{\mathrm{conv}}
    - q_z.
\end{equation}
%
Here $\rho$, $k$, and $\cpeff$ are the local effective density, thermal
conductivity, and effective heat capacity (which folds in latent heat);
$q_{\mathrm{conv}}$ is convective surface loss; and $q_z$ is the
depth-direction correction (Section~\ref{sec:depth}).
Outside $\partmask$, the same equation holds with pure powder-bed properties
and $\Qrf=0$.

\subsection{Time Integration}

Time integration uses \textbf{explicit forward Euler}:
%
\begin{equation}
  T^{n+1} = T^n + \Delta t\left.\frac{\partial T}{\partial t}\right|^n.
\end{equation}
%
The timestep $\Delta t$ is fixed (default 0.5~s, set via \texttt{thermal.dt\_s}).
No CFL stability constraint is enforced analytically; instead the temperature
increment per step is hard-capped:
%
\begin{equation}
  \Delta T_{\mathrm{step}} = \mathrm{clip}\!\left(\Delta t\cdot\dot{T},\;
    -\Delta T_{\max},\;+\Delta T_{\max}\right),
\end{equation}
%
with default $\Delta T_{\max}=3$~\textdegree C/step.
Temperature is also bounded globally to $[-50,\,450]$~\textdegree C.
Total simulated time is $n_{\mathrm{steps}}\times\Delta t$.

\subsection{Thermal Diffusion}

The diffusion term $\nabla\cdot(k\,\nabla T)$ is discretized as a
second-order central-difference flux divergence.
At each cell face, thermal conductivity is \emph{arithmetically averaged}:
%
\begin{equation}
  k_{i+\nicefrac{1}{2},j} = \tfrac{1}{2}(k_{ij}+k_{i+1,j}),\qquad
  k_{i,j+\nicefrac{1}{2}} = \tfrac{1}{2}(k_{ij}+k_{i,j+1}).
\end{equation}
%
Face heat fluxes are:
%
\begin{equation}
  q^x_{i+\nicefrac{1}{2}} = -k_{i+\nicefrac{1}{2},j}\,
    \frac{T_{i+1,j}-T_{i,j}}{\Delta x},\qquad
  q^y_{j+\nicefrac{1}{2}} = -k_{i,j+\nicefrac{1}{2}}\,
    \frac{T_{i,j+1}-T_{i,j}}{\Delta y},
\end{equation}
%
and the divergence at node $(i,j)$ is:
%
\begin{equation}
  \nabla\cdot(k\,\nabla T)\big|_{ij}
    = \frac{q^x_{i+\nicefrac{1}{2}}-q^x_{i-\nicefrac{1}{2}}}{\Delta x}
    + \frac{q^y_{j+\nicefrac{1}{2}}-q^y_{j-\nicefrac{1}{2}}}{\Delta y}.
\end{equation}

\subsection{Convective Surface Loss}

Convective loss is applied as a sink term at all cells on exposed chamber
boundaries:
%
\begin{equation}
  q_{\mathrm{conv}} = h\,A_{\mathrm{expo}}\,(T-T_\infty),
\end{equation}
%
where $A_{\mathrm{expo}}$~[m$^{-1}$] is the exposed-boundary area per unit
volume (i.e.\ $1/\Delta x$ or $1/\Delta y$ for cells on a given edge),
$T_\infty$ is the ambient temperature, and $h$~[W/m$^2$K] is the heat transfer
coefficient.
Two models are available:
\begin{itemize}[leftmargin=2em, itemsep=2pt]
  \item \textbf{Constant-$h$:} user-specified via \texttt{thermal.convection\_h\_w\_per\_m2k}.
  \item \textbf{Natural chimney} (parallel-plate Rayleigh correlation):
    $h$ is computed locally from
    $\mathrm{Ra} = L^3 g\beta\rho^2 c_p\Delta T/(k_{\mathrm{air}}\mu)$
    using Sutherland viscosity and $k_{\mathrm{air}}\propto T^{0.9}$,
    giving $h = k_{\mathrm{air}}\,\mathrm{Ra}/(24H)$, consistent with
    the COMSOL narrow-chimney parallel-plate formulation.
\end{itemize}

\subsection{Depth-Direction Correction}
\label{sec:depth}

Because the simulation is 2D, heat loss through the powder bed above and below
the part (in $z$) is approximated via a series thermal-resistance slab model:
%
\begin{equation}
  q_z = k\,(T-T_\infty)\,\frac{\pi^2}{4\,L_{\mathrm{eff}}^2}
  \bigg|_{\partmask},
\end{equation}
%
where
$L_{\mathrm{eff}} = L_p/2 + L_{\mathrm{pow}}\,(k_{\mathrm{pow}}/k)$
combines the part half-depth $L_p/2$ and the powder-layer half-thickness
$L_{\mathrm{pow}}=L_c/2-L_p/2$ weighted by the conductivity ratio.
This approximates the first-mode solution of 1D slab conduction in $z$.

% ============================================================
\section{Material Properties}

\subsection{Property Mixing}

All material properties inside $\partmask$ are blended between solid-state and
liquid-state values using the local melt fraction $\phimelt\in[0,1]$ and the
relative density $\rhorel\in[0,1]$:
%
\begin{align}
  \rho_{\mathrm{s,eff}} &= \rho_{\mathrm{pow}} + \rhorel(\rho_s - \rho_{\mathrm{pow}}),\\
  k_{\mathrm{s,eff}}    &= k_{\mathrm{pow}}   + \rhorel(k_s   - k_{\mathrm{pow}}),\\
  \rho_{\mathrm{local}} &= (1-\phimelt)\,\rho_{\mathrm{s,eff}} + \phimelt\,\rho_\ell,\\
  k_{\mathrm{local}}    &= (1-\phimelt)\,k_{\mathrm{s,eff}}   + \phimelt\,k_\ell,\\
  c_{p,\mathrm{base}}   &= (1-\phimelt)\,c_{p,s}              + \phimelt\,c_{p,\ell}.
\end{align}
%
Outside the part, $\rho=\rho_{\mathrm{pow}}$, $k=k_{\mathrm{pow}}$,
$c_p=c_{p,\mathrm{pow}}$ (virgin powder bed).

\subsection{RF Conductivity}

The RF conductivity field used to build $\gamma$ blends between virgin powder
and doped-part conductivities via the subpixel fill fraction:
%
\begin{equation}
  \sigma_{ij} = \sigma_{\mathrm{virgin}} + f_{ij}(\sigma_{\mathrm{doped}}-\sigma_{\mathrm{virgin}}).
\end{equation}
%
An optional temperature and density dependence can be applied:
%
\begin{equation}
  \sigma(T,\rhorel) = \sigma_0\,(1+\alpha_T(T-T_{\mathrm{ref}}))\,(1+\alpha_\rho\,\rhorel),
\end{equation}
%
where $\alpha_T$ and $\alpha_\rho$ are the temperature and density coefficients
(both default to zero).

% ============================================================
\section{Phase Change and Latent Heat}

The melt fraction $\phimelt$ and its temperature derivative $d\phimelt/dT$
enter $\cpeff$ through:
%
\begin{equation}
  \cpeff = c_{p,\mathrm{base}} + L\,\frac{d\phimelt}{dT},
\end{equation}
%
so latent heat is captured implicitly via the effective heat capacity
rather than as a separate enthalpy source term.
Two phase models are available.

\subsection{Baseline: Smoothed Heaviside}

A linear ramp between solidus $T_s=T_{\mathrm{pc}}-\Delta T_{\mathrm{pc}}/2$
and liquidus $T_\ell=T_{\mathrm{pc}}+\Delta T_{\mathrm{pc}}/2$:
%
\begin{equation}
  \phimelt =
  \begin{cases}
    0 & T < T_s,\\
    (T-T_s)/\Delta T_{\mathrm{pc}} & T_s\le T\le T_\ell,\\
    1 & T > T_\ell,
  \end{cases}
  \qquad
  \frac{d\phimelt}{dT} = \frac{1}{\Delta T_{\mathrm{pc}}}\quad\text{(in window)}.
\end{equation}

\subsection{DSC-Calibrated Apparent Heat Capacity}

A triangular/Gaussian-blended profile is fitted to differential scanning
calorimetry (DSC) measurements.
For PA12 the calibrated values are:
%
\begin{equation}
  T_{\mathrm{onset}} = 171\,\mathrm{^\circ C},\quad
  T_{\mathrm{peak}}  = 180.8\,\mathrm{^\circ C},\quad
  T_{\mathrm{end}}   = 186\,\mathrm{^\circ C},\quad
  L = 101.7\;\mathrm{kJ/kg}.
\end{equation}
%
The curve $d\phimelt/dT$ is evaluated directly from the DSC profile and
fed into $\cpeff$, allowing precise representation of the asymmetric
melt transition observed in semicrystalline polymers.

% ============================================================
\section{Densification Model}
\label{sec:dens}

Relative density $\rhorel$ evolves according to
$d\rhorel/dt = R(T,\phimelt,\rhorel)$,
with $\rhorel$ bounded to $[\rho_{\mathrm{pow}}/\rho_s,\,1]$.
Three rate models are available.

\subsection{Dual-Mechanism (\texttt{physics\_dual})}

Two kinetic mechanisms operate additively.

\paragraph{Solid-state (diffusion/creep):}
%
\begin{equation}
  k_{\mathrm{ss}} = k_{0,\mathrm{ss}}\exp\!\left(-\frac{E_{a,\mathrm{ss}}}{RT}\right),
  \qquad
  R_{\mathrm{ss}} = k_{\mathrm{ss}}\,(1-\phimelt)^{n_s}\,\rho_{\mathrm{term}}.
\end{equation}

\paragraph{Liquid-phase (viscous-capillary):}
%
\begin{equation}
  k_{\mathrm{liq}} = \frac{C\,\gamma_{\mathrm{surf}}}{\eta(T)\,r_p},
  \qquad
  R_{\mathrm{liq}} = k_{\mathrm{liq}}\,
    \left(\frac{\phimelt-\phi_{\mathrm{thr}}}{1-\phi_{\mathrm{thr}}}\right)^{n_\ell}
    \rho_{\mathrm{term}},
\end{equation}
%
where $\eta(T)=\eta_{\mathrm{ref}}\exp[E_{a,\eta}/R\,(1/T-1/T_{\mathrm{ref}})]$
is the Arrhenius melt viscosity, $\gamma_{\mathrm{surf}}$ is surface tension,
$r_p$ is particle radius, and $C$ is a geometric packing factor.
The net densification rate is $R = R_{\mathrm{ss}} + R_{\mathrm{liq}}$.

\subsection{Experimental PA12 Model (\texttt{experimental\_viscous\_capillary\_pa12})}

Single viscous-capillary mechanism calibrated to PA12.
Two viscosity sub-models are supported:

\paragraph{Arrhenius:}
$\eta(T) = \eta_{\mathrm{ref}}\exp\!\left[E_{a,\eta}/R\,(1/T-1/T_{\mathrm{ref}})\right]$.

\paragraph{Williams-Landel-Ferry (WLF):}
%
\begin{equation}
  \log_{10}\eta = \log_{10}\eta_{\mathrm{ref}}
    - \frac{C_1(T-T_{\mathrm{ref}})}{C_2+(T-T_{\mathrm{ref}})}.
\end{equation}
%
The rate includes a melt-activation drive
$\phi^*=\bigl((\phimelt-\phi_{\mathrm{thr}})/(1-\phi_{\mathrm{thr}})\bigr)^p$
and a porosity drive $(1-\rhorel)^q$:
%
\begin{equation}
  R = \frac{C\,\gamma_{\mathrm{surf}}}{\eta(T)\,r_p}\cdot(\phi^*)^p\cdot(1-\rhorel)^q.
\end{equation}

\subsection{Simple Viscous-Capillary (\texttt{viscous\_capillary})}

Single Arrhenius-viscosity capillary model; intermediate in complexity
between the baseline and PA12-calibrated variants.

% ============================================================
\section{Coupled Time-Integration Loop}

The three physical fields --- electric potential, temperature, and relative
density --- are advanced by a split-operator explicit loop:

\medskip
\noindent
\begin{tabular}{ll}
\textbf{for} $n = 0,1,\ldots,N_{\mathrm{steps}}$: & \\
\quad\textbf{if} $n \bmod N_{\mathrm{update}}=0$: &
  solve EQS~\eqref{eq:eqs} $\to \Qrf^n$ \quad(sparse LU)\\
\quad $T^{n+1} = T^n + \Delta t\,\mathcal{L}_T(T^n,\Qrf^n)$ &
  (explicit FD, vectorized)\\
\quad $\rhorel^{n+1} = \rhorel^n + \Delta t\,R(T^{n+1},\phimelt^{n+1},\rhorel^n)$ &
  (explicit ODE, vectorized)
\end{tabular}

\medskip
\noindent
where $\mathcal{L}_T$ is the full right-hand side of Equation~\eqref{eq:heat}
and $N_{\mathrm{update}}$ is the EQS re-solve interval
(default: every step; configurable to re-solve only at rotation events
for turntable runs).

\paragraph{Computational cost.}
The EQS sparse direct solve dominates (milliseconds per solve on a
$120\times120$ grid).
The thermal and densification updates are vectorized NumPy operations on the
same grid and are comparatively inexpensive.
Total wall time for a 6-minute exposure at default grid resolution is
typically 30--120~s depending on the EQS update interval and densification
model selected.

% ============================================================
\section{Summary of Key Parameters}

\begin{center}
\renewcommand{\arraystretch}{1.3}
\begin{tabular}{lllp{5.2cm}}
\toprule
\textbf{Parameter} & \textbf{Symbol} & \textbf{Default} & \textbf{Config key} \\
\midrule
Grid size         & $N_x\times N_y$        & $120\times120$   & \texttt{geometry.grid\_nx/ny} \\
Chamber size      & $L_x\times L_y$        & $60\times60$~mm  & \texttt{geometry.chamber\_x/y} \\
Timestep          & $\Delta t$             & $0.5$~s          & \texttt{thermal.dt\_s} \\
Max $\Delta T$/step & $\Delta T_{\max}$    & $3$~\textdegree C  & \texttt{thermal.max\_deltaT\_per\_step\_c} \\
RF frequency      & $f$                    & $27.12$~MHz      & \texttt{electric.frequency\_hz} \\
EQS update interval & $N_{\mathrm{update}}$ & 1 step          & \texttt{electric.update\_interval} \\
DSC melt onset    & $T_{\mathrm{onset}}$   & $171$~\textdegree C   & \texttt{dsc\_profile.melt\_onset\_c} \\
DSC melt peak     & $T_{\mathrm{peak}}$    & $180.8$~\textdegree C & \texttt{dsc\_profile.melt\_peak\_c} \\
DSC melt end      & $T_{\mathrm{end}}$     & $186$~\textdegree C   & \texttt{dsc\_profile.melt\_end\_c} \\
Latent heat       & $L$                    & $101.7$~kJ/kg    & \texttt{dsc\_profile.lat\_heat\_j\_per\_kg} \\
\bottomrule
\end{tabular}
\end{center}

\end{document}
